% Lösungen Geraden · Punktprobe und Punkte auf der Geraden · Prüfungs-Fokus Abitur GK (zusammenbau v0.6)
\einheitenkopf{Einheit 2 · Punktprobe und Punkte auf der Geraden}

\erg{1}{a) aus der zweiten oder der dritten – die erste Komponente des Richtungsvektors ist null \quad b) Aus der ersten Koordinate $r = 3$; zweite: $2 + 3 = 5$ ✓, dritte: $-1 + 9 = 8$ ✓ – $P$ liegt auf $g$.}
\erg{2}{a) Jeder Punkt von $g$ hat die zweite Koordinate $-1$ (Richtungskomponente null), $A$ hat 3. \quad b) Jeder Punkt von $h$ hat die dritte Koordinate 6 (Richtungskomponente null), $B$ hat 5. \quad c) Jeder Punkt von $k$ hat die erste Koordinate $-3$ (Richtungskomponente null), $C$ hat 1 – dass die zweite und die dritte Koordinate zu $t = 3$ passen, ändert daran nichts. \quad d) $1 - r = -2$ liefert $r = 3$; $y_B = 3 + 6 = 9$, $z_B = 2 - 9 = -7$ \quad e) $-1 + t = 2$ liefert $t = 3$; $x_P = 0 + 6 = 6$, $z_P = 4 - 6 = -2$ \quad f) $0 + t = -3$ liefert $t = -3$; $x_Q = 3 + 6 = 9$, $y_Q = -2 - 12 = -14$ \quad g) $|\vec u| = 3$, also $A + 3 \cdot \vec u = (9 | -2 | 8)$ (oder $A - 3 \cdot \vec u = (-3 | 4 | -4)$) \quad h) $|\vec u| = 5$, also $A + 2 \cdot \vec u = (8 | 4 | -7)$ (oder $A - 2 \cdot \vec u = (-8 | 4 | 5)$) \quad i) $|\vec u| = 3$, Faktor $4{,}5 : 3 = 1{,}5$, also $A + 1{,}5 \cdot \vec u = (-0{,}5 | 4 | 6)$ \quad j) Ansatz $(2 | 6 | 0) = (8 | 0 | 0) + t \cdot (-8 | 8 | 0)$: erste und zweite Koordinate liefern $t = 0{,}75$, die dritte stimmt, und $0 \le 0{,}75 \le 1$ – $T$ liegt auf der Strecke. \quad k) Ansatz $(-3 | 7 | 0) = (6 | -2 | 0) + t \cdot (-4 | 4 | 0)$: erste und zweite Koordinate liefern $t = 2{,}25$, die dritte stimmt – $T$ liegt auf der Geraden $EF$, wegen $2{,}25 > 1$ aber nicht auf dem Zaun. \quad l) Ansatz $L = A + t \cdot (6 | 3 | -3)$: aus der ersten Koordinate $t = \tfrac{2}{3}$; zweite: $1 + 2 = 3$ ✓, dritte: $4 - 2 = 2$ ✓; $0 \le \tfrac{2}{3} \le 1$ – $L$ liegt auf dem Seil.}
\erg{3}{a) $\overrightarrow{AB} = (2 | 1 | -2)$, $\overrightarrow{AC} = (6 | 3 | -6) = 3 \cdot \overrightarrow{AB}$ – die Verbindungsvektoren sind kollinear, $A$, $B$, $C$ liegen auf einer Geraden. \quad b) $\overrightarrow{AB} = (1 | -2 | 1)$, $\overrightarrow{AC} = (-2 | 4 | -2) = -2 \cdot \overrightarrow{AB}$ – kollinear, die drei Punkte liegen auf einer Geraden. \quad c) $\overrightarrow{AC} = (6 | 3 | 10)$, aber $3 \cdot \overrightarrow{AB} = (6 | 3 | 9)$ – kein Vielfaches, die Punkte liegen nicht auf einer Geraden. \quad d) $|\overrightarrow{AE}| = 21$ LE, 84 m sind $16{,}8$ LE, also $\lambda = 16{,}8 : 21 = 0{,}8$; $x_3 = 2 + 0{,}8 \cdot 9 = 9{,}2$ LE, Höhe: $9{,}2 \cdot 5 = 46$ m \quad e) $|\overrightarrow{AE}| = 12$ LE, 30 m sind 3 LE, also $\lambda = 0{,}25$; $x_3 = 4 + 0{,}25 \cdot 8 = 6$ LE, Höhe: $6 \cdot 10 = 60$ m \quad f) $|\overrightarrow{AE}| = 12$ LE, 180 m sind 9 LE, also $\lambda = 0{,}75$; $x_3 = 0 + 0{,}75 \cdot 8 = 6$ LE, Höhe: $6 \cdot 20 = 120$ m \quad g) $\vec x = (5 | 0 | 0) + r \cdot (-4 | 4 | 6)$; aus der ersten Koordinate $r = 0{,}5$, zweite: $0 + 2 = 2$ ✓, dritte: $0 + 3 = 3 \ne 5$ – $S$ liegt nicht auf der Geraden. \quad h) $\vec x = (0 | 6 | 1) + r \cdot (4 | -6 | 3)$; aus der ersten Koordinate $r = 0{,}5$, zweite: $6 - 3 = 3$ ✓, dritte: $1 + 1{,}5 = 2{,}5 \ne 2$ – $S$ liegt nicht auf der Geraden. \quad i) $\vec x = (2 | 0 | 0) + r \cdot (-4 | 6 | 6)$; aus der ersten Koordinate $r = 0{,}5$, zweite: $0 + 3 = 3$ ✓, dritte: $0 + 3 = 3$ ✓ – $S$ liegt auf der Geraden (in der Mitte zwischen $A$ und $G$).}
